\section{Branches, opened atoms, and lossless cells}
\label{sec:literal-frame}

\subsection{The resolved branch is the primary key}

Fix a soluble nonempty canonical branch
\[
 \beta=(\theta,c,\kappa),\qquad
 \theta=(L,\gamma,\ell,j,\sigma,v,\iota).
\]
Put
\begin{equation}
 b_\beta=c_\beta\kappa_\beta,\qquad
 g_\beta=(b_\beta,\sigma_\beta),\qquad
 B_\beta=\frac{b_\beta}{g_\beta}.
\label{eq:branch-content}
\end{equation}
Content resolution gives one progression
\[
 t_\beta(z)=\tau_\beta+B_\beta z,
 \qquad z\in I_\beta^\ast.
\]
Its cardinality is \(N_\beta=\#I_\beta^\ast\).  On the
nonconstant resonant carrier, the TPC-99 exact resonance-window
lemma gives \(N_\beta\le q_X\) for every nonzero resonant
frequency, and defines the width
\[
 H_q(N)=
 \min\!\left\{\frac{q-1}{2},\left\lfloor\frac qN\right\rfloor\right\},
\]
is
\begin{equation}
 \mathsf H_\beta=H_{q_X}(N_\beta).
\label{eq:branch-width}
\end{equation}
Branches with \(N_\beta>q_X\) have no nonzero resonance and are not
entered into a width census.  The source-level
\(q_X\mid u_p\) bound below does not use \(\mathsf H_\beta\) and may be
applied before this restriction.
\(\mathsf H_\beta\) denotes only this resonance width; it is not the
robustness-content variable denoted \(H_\beta\) in the general
TPC-94 normalization.  On the canonical branch that earlier content
equals \(1\).
The phase slope and native multiplier are constant on this
\(z\)-fiber:
\begin{equation}
 \Omega_\beta
 =
 \ell_\beta v_\beta\sigma_\beta B_\beta,
\qquad
 \omega_\beta
 =
 \varepsilon_\beta\frac{\Omega_\beta}{c_\beta}
 \pmod{q_X}.
\label{eq:branch-multiplier}
\end{equation}

The actual source values are not constant on the fiber.  For
\(z\in I_\beta^\ast\), define
\begin{equation}
 p=(\beta,z),\qquad
 u_p=U_\theta(t_\beta(z)),\qquad
 m_p=M_\theta(t_\beta(z)).
\label{eq:opened-atom-values}
\end{equation}
Here \(p\) is an \(r\)-free opened atom.  Since
\[
 M_\theta(t)=\ell_\theta v_\theta D_\theta(t),
\qquad
 \sigma_\theta U_\theta(t)
 =
 M_\theta(t)j_\theta+h_0,
\]
every atom satisfies the exact source identity
\begin{equation}
 m_pj_p+h_0=\sigma_pu_p,
\qquad
 j_p=j_\theta,\quad\sigma_p=\sigma_\theta.
\label{eq:atom-source-identity}
\end{equation}
Different \(z\)'s of one branch generally give different \(u_p\)
and \(m_p\), while \(j_p,\sigma_p,\omega_\beta,\mathsf H_\beta\) remain
fixed.

\subsection{Exact atomization of the TPC-99 weight}

Let \(A_{\beta,X}\) be the literal outer branch coefficient from
TPC-94 and let \(\mathcal W_{\beta,X}(z)\) contain the actual smooth,
periodic, extraction, and coprimality masks.  Define
\begin{align}
 w_{\beta,z}
 ={}&
 |A_{\beta,X}|\,
 |\mathcal W_{\beta,X}(z)|
 \mu^2(\widetilde D_\beta(z))
 \mu^2(V_\beta(z)),
\label{eq:atom-weight}\\
 w_\beta
 ={}&
 \sum_{z\in I_\beta^\ast}w_{\beta,z}.
\label{eq:branch-weight}
\end{align}
Thus \(w_\beta\) is exactly the proof-carrying TPC-99 branch weight:
in the notation of that paper,
\[
 w_\beta=w_{\beta,X}
 =|A_{\beta,X}|\,\mathcal A_{\beta,X}.
\]
The refinement \eqref{eq:branch-weight} is an equality, not a second
copy of the branch mass.

For exact width \(H\), the native census is therefore
\begin{align}
 a_H(\omega)
 &=
 \sum_{\substack{\beta:
       \mathsf H_\beta=H\\\omega_\beta=\omega}}w_\beta
\notag\\
 &=
 \sum_{\substack{p=(\beta,z):
       \mathsf H_\beta=H\\\omega_\beta=\omega}}w_{\beta,z}.
\label{eq:atomized-census}
\end{align}
The frequency \(r\) is not part of \(p\); it remains in the
low-window weight \(b_X(r)=|\mathfrak m_{K,X}(r)|\).

\begin{remark}[Width is not cell occupancy]
\label{rem:width-not-cell}
Throughout the paper,
\[
 \mathsf H_p=\mathsf H_\beta=H_{q_X}(\#I_\beta^\ast).
\]
It is the resonance width of the original branch \(z\)-fiber.  It
is never replaced by \(|J_C|\), by the number of opened atoms in a
cell, or by the size of either ultra-divisor submass.
\end{remark}

\subsection{The augmented atom record}

TPC-93 gives an explicit inverse from a child \((\theta,t)\) to the
preprojector source atom.  Augment \(p=(\beta,z)\) by:
\begin{itemize}[leftmargin=2em]
\item the child coordinate
  \(t_p=\tau_\beta+B_\beta z\) and the inverse source atom, including
  polarization, actual moving and opposite rows, \(j_p\), and the
  opened value \(u_p\);
\item the projector divisor \(v_p\), interval component, actual
  masks, prefix, residue, smooth, and reconstruction labels;
\item \(c_\beta,\kappa_\beta,b_\beta,g_\beta,B_\beta\), the branch
  interval \(I_\beta^\ast\), \(N_\beta,\mathsf H_\beta\), and \(z\); and
\item every label needed to reverse the exact-content and source
  reconstructions; the atom weight \(w_{\beta,z}\) is attached as a
  coefficient payload, not used as a grouping label.
\end{itemize}
Redundant derived values such as \(m_p,u_p,\sigma_p\) are retained
on purpose; they are provenance checks, not new multiplicities.

\begin{definition}[Fine provenance cell]
\label{def:fine-cell}
Flatten the augmented provenance record: \(\beta\) is not kept as an
opaque token that could secretly retain \(j_\theta\).  Let
\(\operatorname{del}_j(p)\) delete the single native coordinate
\(j_p\) and its redundant copies, while deleting no independent
coordinate.  The attached numerical weight is not part of this
projection.  A fine cell \(C\) is one fiber of
\(\operatorname{del}_j\).  Write
\[
 P_C=\{p:\operatorname{del}_j(p)=C\},\qquad
 J_C=\{j_p:p\in P_C\}.
\]
\end{definition}

\begin{theorem}[Lossless opened-atom disintegration]
\label{thm:lossless-atom-cells}
Assume the TPC-93 child-to-source inverse and the TPC-94 exact
content reconstruction.  Then:
\begin{enumerate}[label=(\roman*),leftmargin=2.2em]
\item \((C,j_p)\) reconstructs the unique pair \((\beta,z)\), its
  \(r\)-free source atom, and its weight \(w_{\beta,z}\);
\item in fact \(C\) already determines
  \[
   j_p=\frac{\sigma_pu_p-h_0}{m_p},
  \]
  so every nonempty fine cell \(P_C\) is a singleton;
\item all \(p\in P_C\) have common values
  \(m_C,u_C,\ell_C,v_C,c_C,B_C,\varepsilon_C,\mathsf H_C\); and
\item summing the cells gives exactly the second line of
  \eqref{eq:atomized-census}.
\end{enumerate}
\end{theorem}

\begin{proof}
The flattened augmented record contains every coordinate of
\((\beta,z)\), its TPC-93 inverse source record, and every
reconstruction label.  It contains no opaque branch identifier.
The map \(\operatorname{del}_j\) removes only the native \(j\)
coordinate and direct stored copies, so restoring that coordinate
restores the original record and its attached coefficient.  This
proves (i).  The retained physical coordinates
\(m_p,u_p,\sigma_p\) remain.  Their source identity
\eqref{eq:atom-source-identity} determines the displayed integer
\(j_p\), proving (ii).  Every quantity in (iii) remains in the
deletion record, either as a native coordinate or an explicitly
retained source value.  Finally, the cells form a disjoint partition
of the opened atoms, so (iv) is just finite reindexing of
\eqref{eq:atomized-census}.
\end{proof}

\begin{remark}[No claim that \(\theta\) fixes \(u\)]
The cell fixes the \emph{opened source value}
\(u_C=U_\theta(t_p)\) only after \(z\) has been opened and the source
inverse recorded.  The branch \(\beta\), and in particular
\(\theta\), does not fix \(u\) or \(m\) along its full resolved
fiber.
\end{remark}

\begin{remark}[The fine-cell injection is deliberately vacuous]
\label{rem:singleton-cells}
The singleton conclusion is not an equidistribution statement.  It
is the unavoidable price of retaining the complete source inverse
while deleting only the displayed \(j\)-coordinate: the determinant
identity algebraically recovers that coordinate.  Consequently all
collisions between distinct opened atoms, including same-\(\beta\)
different-\(z\) collisions, occur in the cross-cell ledger.
\end{remark}

\subsection{Strong no-wrap consequences}

The strengthened prime of TPC-98 may be taken, still with
\(q_X\asymp Q_X\), larger than every actual source row, every
relevant row difference, and the orbit diameter.  On active
canonical support,
\begin{equation}
 0<m_p,\ell_p,v_p,c_\beta,B_\beta,b_\beta<q_X,
\qquad
 \operatorname{diam}(\mathcal J_X)<q_X.
\label{eq:atom-no-wrap}
\end{equation}
Hence \(m_p,\ell_p,v_p,c_\beta,B_\beta\) are invertible modulo
\(q_X\).  On the nonconstant branch,
\begin{equation}
 q_X\nmid\Omega_\beta
 \quad\Longleftrightarrow\quad
 q_X\nmid\sigma_\beta,
\qquad
 \omega_\beta\in\mathbb F_{q_X}^{\times}.
\label{eq:atom-nonconstant}
\end{equation}
The equivalence uses the literal bound \(B_\beta<q_X\) and the
nonvanishing of \(\ell_\beta,v_\beta\) modulo \(q_X\)
\citep{WangTPC98}.
